% Source: sources/returns/S6_04_MAPS_MANIFOLDS/sections/manifolds.tex:3025–3270
\paragraph{Theorem D.2.15 (Proved: all-order $S^2\times S^4$ outputs).}
Take the OS56 type-$II$ profile with ordered elliptic fibres
\[
                         I_2,\ II,\ I_5,\ I_3.
\]
Let $G$ be its marked Mordell--Weil generator, take
$P=30G$, $Q=G$, and put the degree-one linearization zero
at $I_2$. Use the two-component toroidal filling there,
the original primitive one-node fillings at $I_5,I_3$,
and the logarithmic filling at $II$ with
\[
 m=6k,\qquad v=n\gamma+r\delta,\qquad
 k>0,\quad\gcd(n,6k)=1,\quad r\in\mathbb Z.
\]
The resulting compact smooth threefold has
\[
\begin{aligned}
 \pi_1(X)&=H_1(X)=\mathbb Z/|n|,\qquad H_5(X)=0,\\
 H_2(X)&=H_4(X)=\mathbb Z\oplus\mathbb Z/|n|,
 \qquad H_3(X)=\mathbb Z/|n|,\\
 H_0(X)&=H_6(X)=\mathbb Z,\qquad \chi(X)=4.
\end{aligned}
\]
Every simply connected member, equivalently $n=\pm1$,
is oriented diffeomorphic to $S^2\times S^4$.
In particular this fixed configuration supplies that
smooth type for every logarithmic multiplicity $6k$.
\begin{proof}
\emph{The actual local geometries.}
The generator has height $1/30$, and $P$ is narrow at
every fibre. The component residues of $Q$ are $1$ in
each of $\mathbb Z/2$, $\mathbb Z/5$, $\mathbb Z/3$,
and zero at $II$. The mixed height is one, so the same
vertical-divisor argument as in D.2.12 gives
$M\otimes t_Q^*M^{-1}=\pi^*\mathcal O(1)$.
Its section can have its simple zero on the whole $I_2$
fibre. The two weight-zero multiplicative fibres give
the primitive one-node models proved in D.2.12.
The type-$II$ good-reduction model has linear order six,
with invariant basis $\delta,\gamma$; its higher-order
free affine quotients are supplied by the same holomorphic
construction as in D.2.12.

At the zero of the linearization, the Tate parameter has
order two. Because $P$ is narrow, its line has order zero
in the angular Tate direction, while $Q$ has component
residue one. The two valuation periods of the torus
uniformization are therefore
\[
                 (2,0),\qquad(1,1).
 \tag{D.I1}
\]
The second coordinate one is the simple zero of the
linearization. Holomorphic units do not change this
integral lattice. Use the standard unit-triangle Mumford
fan and quotient by the lattice $\Gamma$ in (D.I1).
It has index two, so the central fibre has two normalized
$dP_6$ components, six double curves and four triple
points. Its total space is smooth: all the fan cones are
unimodular, and the contracting period action is free
and properly discontinuous on the chosen tube.
This is the specified two-component filling, not a
primitive one-component replacement of its monodromy.

\emph{The integral cusp defect.}
Let $U$ be a compact tubular neighborhood of this cusp,
and $M=\partial U$. The fan symmetry
\[
 F=\begin{pmatrix}1&1\\0&-1\end{pmatrix},\qquad
 F\begin{pmatrix}2&1\\0&1\end{pmatrix}
     =\begin{pmatrix}2&2\\0&-1\end{pmatrix}
\]
preserves the triangular fan and identifies this lattice
with $2\mathbb Z\oplus\mathbb Z$.
Consequently the integral normalization resolution in
D.1.2, at $d=2$, applies directly and gives
\[
 \pi_1(U)=\mathbb Z^2,\qquad
 (H_0,H_1,H_2,H_3,H_4)(U)
   =(\mathbb Z,\mathbb Z^2,\mathbb Z^5,\mathbb Z^2,\mathbb Z^2).
 \tag{D.I2}
\]
This local normalization calculation does not require the
global covering degree to be coprime to six.

The cusp monodromy $A_3$ has $\operatorname{rank}(A_3-I)=2$
and Smith entries $(1,2)$, while
$\bigwedge^2A_3-I$ has two unit Smith entries.
The Wang sequences therefore give
\[
 H_1(M)=\mathbb Z^3\oplus\mathbb Z/2,
 \qquad H_2(M)=\mathbb Z^6.
\]
The cusp kills the full saturated angular lattice
$\ker(A_3-I)$, and its canonical meridian bounds the
extended section disc. Hence the map on $H_1$ is onto
$H_1(U)=\mathbb Z^2$, with kernel
$\mathbb Z\oplus\mathbb Z/2$.
Poincare--Lefschetz duality gives
$H_2(U,M)=H^4(U)=\mathbb Z^2$ and $H_1(U,M)=0$.
The relative sequence then shows that
\[
 H_2(U)/\operatorname{im}H_2(M)\cong\mathbb Z.
 \tag{D.I3}
\]
In particular this quotient has no hidden two-torsion.
The two angular normal tori bound in the toric model.
They project to the full primitive fixed lattice in the
Wang sequence for $H_2(M)$, so they form a primitive
rank-two subgroup. By (D.I2)--(D.I3) the kernel of
$H_2(M)\to H_2(U)$ has exactly rank two; these are
therefore its whole kernel.
Moreover the remaining image is exactly the regular-fibre
specialization image. Its cokernel is the free group
in (D.I3).

\emph{The global integral map.}
Cyclically reorder the original four monodromies as
$A_0,A_1,A_2,A_3$ corresponding to $II,I_5,I_3,I_2$.
The full matrices are retained in the accompanying certificate
\texttt{manifolds\_i2\_cusp.json}.
They satisfy $A_0\cdots A_3=I$, $A_0^6=I$ and
fix $\delta$ and the covector $\gamma^*$.
Here $A_0\gamma=\gamma$ and
\[
 \omega=30e_1^*e_2^*-\delta^*\gamma^*
\]
is the primitive exterior coinvariant detector.
The exterior presentation has five unit Smith entries,
so its group is free of rank one.

There is one integral issue in degree one that cannot be
discarded. The combined matrix $[A_i-I]_{i=0}^3$ has
Smith entries $(1,1,2)$, so the regular family's fibre
coinvariants are $\mathbb Z\oplus\mathbb Z/2$.
The torsion in the cusp boundary is represented by
$-\delta$, and its image in this global $\mathbb Z/2$
has coefficient $-1$. Thus the torsion summands in
the degree-one Mayer--Vietoris map are mapped
isomorphically. Replacing the cusp image by its actual
saturated angular vanishing lattice gives Smith entries
$(1,1,1)$; it kills the entire primitive
$\ker\gamma^*$ in van Kampen. The base relation kills
the remaining finite meridian, leaving just
$n\gamma=0$. This proves $\pi_1(X)=\mathbb Z/|n|$.

Use the boundary maps $j_i$ of D.2.12 and put
\[
 K=[j_1\ker(A_1-I)\ j_2\ker(A_2-I)
       \ j_3\ker(A_3-I)\ j_0(\gamma)].
\]
The exact integer certificate gives a left inverse
$LK=I_9$, and $(Lj_0\delta)_9=0$.
As before, these actual fixed-loop sweeps have zero
$\omega$-area. After the six one-node normal tori and
the two cusp normal tori, the regular-family $H_2$ is
therefore $\mathbb Z\langle s,t\rangle$ with
$t=j_0\gamma$.
The finite vertical coinvariants are generated by
$\delta$, and
$\mathcal N_m\delta=m\delta$,
$\mathcal N_m\gamma=m\gamma$.
Writing $m\alpha+n\beta=1$, its two normal tori give
\[
                        \begin{pmatrix}n&-r\alpha\\0&1\end{pmatrix}.
 \tag{D.I4}
\]
Their cokernel is $\mathbb Z/|n|$ for all parameters.
All these equalities are integral vector identities in
the retained matrices.

The degree-two boundary maps of the other fillings are
onto. At the cusp their cokernel is the free group
(D.I3). Thus the full degree-two Mayer--Vietoris cokernel
is $\mathbb Z\oplus\mathbb Z/|n|$; the extension by
the free cusp defect splits. The degree-one source and
target both have free rank fourteen and one copy of
$\mathbb Z/2$. Their torsion map is the isomorphism
just checked, and their total cokernel is the finite
group $H_1(X)$. Their kernel is therefore zero.
This proves the asserted $H_2$. Four toric triple points
give Euler characteristic four; integral duality then
gives all the remaining groups.

\emph{Smooth recognition at $n=\pm1$.}
In this case $H_2(X)=\mathbb Z$, and the cusp map
$H_2(U)\to H_2(X)$ is onto, because (D.I4) kills the
entire regular-family contribution. The local Mumford
threefold has trivial canonical line. Hence $c_1(X)$
evaluates as zero on every $H_2$ class, so $c_1(X)=0$
and $X$ is spin.

The angular $\delta$ action is semifree with two fixed
rational curves. The one-node pieces are free as before.
At the finite logarithmic fibre, its invariant
$\gamma^*$ coordinate changes by $n/m$ under a deck
generator; for $n=\pm1$ this prevents any additional
stabilizer of the $\delta$ action. At the cusp there
are two vertical-edge orbits, each with the normal
weights and degrees (D.S11).

It remains essential to check that these two fixed curves
have zero homology class; specialization is not onto in
(D.I3). Colour a lattice vertex $(a,b)$ by $a+b$ modulo
two to label its normalized component. For a vertical
edge from $(a,b)$ to $(a,b+1)$, the two adjacent third
vertices are $(a+1,b)$ and $(a-1,b+1)$.
Each colour occurs once among the endpoints and once
among the third vertices. The toric wall formula thus
gives degree $-1+1=0$ for the line of either global
component on either fixed curve. Their images in
$H_2(U,M)$ are exactly these component intersection
numbers, so both images are zero. By the relative
sequence both curves lie in the boundary image, which
is precisely the regular-fibre specialization image.
Equation (D.I4), with $n=\pm1$, now kills both classes
in $H_2(X)$.

The equivariant surgery and spin-lift argument of D.2.14
can therefore be performed on both disjoint curves.
Whitney isotopy applies also to their disjoint union,
and their normal framings are unique. It gives
\[
            Y\cong X\#_2(S^3\times S^3)
\]
with a free circle action and simply connected spin
quotient $B$. The homotopy sequence gives
$H_2(B)=\mathbb Z^2$ with primitive Euler class.
The oriented Smale--Barden theorem identifies the pair
with $((S^2\times S^3)\#(S^2\times S^3),(x,0))$.
Applying \cite[Theorem A and Theorem B(2)]{manGR25}
gives
\[
              Y\cong(S^2\times S^4)\#_2(S^3\times S^3).
\]
Zhubr cancellation as in D.2.14 proves
$X\cong S^2\times S^4$ with its oriented smooth type.

The exact program \texttt{code/manifolds\_i2\_cusp.py}
verifies the primitive graph and exterior maps, the
nonzero boundary-torsion image, the actual saturated
vanishing lattice, the two fixed-curve degree vectors,
and 630 full $10\times10$ Smith forms.
The norm identities and (D.I4) prove the result without
any bounds on $k,n,r$. No priority claim is made.
\end{proof}
